\documentclass[10pt,a4paper]{amsart}
\usepackage[margin=19mm]{geometry}
\usepackage{amsmath,amssymb,mathtools}
\usepackage[scr=rsfs,cal=euler]{mathalfa}
\usepackage{xfrac}
\usepackage[unicode,hidelinks,bookmarks=false]{hyperref}
\newcommand{\ie}{i.e.\ }
\newcommand{\cf}{cf.\ }
\newcommand{\resp}{respectively\ }
\newcommand{\Subsection}{\subsection}
\providecommand{\nobreakdash}{\nobreak\hbox{-}}
\setlength{\emergencystretch}{2em}
\multlinegap0pt
\usepackage{amssymb}
\newtheorem{theorem}{Theorem}
\title{Refinements of Alon–Babai–Suzuki-type intersection theorems via non-shadows and binomial support}
\author{Jiangdong Ai, Mingyu Liu}
\date{Source: arXiv:2603.10309v2 (CC BY 4.0)}
\begin{document}
\maketitle

\noindent\emph{Source and licence.} Refinements of Alon–Babai–Suzuki-type intersection theorems via non-shadows and binomial support. Version arXiv:2603.10309v2 (CC BY 4.0), distributed by arXiv under the Creative Commons Attribution 4.0 International licence, http://creativecommons.org/licenses/by/4.0/, as recorded in the arXiv licence metadata for this exact version and shown on the abstract page https://arxiv.org/abs/2603.10309v2. Full licence notice: \texttt{licenses/arxiv-2603.10309-CC-BY-4.0.txt}.\par\smallskip

\noindent\textit{Editorial scope.} Counted unit: the article's theorem theorem (source0.tex lines 357--372). The article's complete original proof of it is delivered with it (source0.tex lines 374--475, one \texttt{proof} environment of 102 lines). No mathematics is written, repaired or re-derived: the statement and the proof are the article's own lines, in the article's own notation, and no retained line is substituted or re-flowed.  No further source-local context is needed: every cross-reference of every retained line resolves inside the two delivered spans.  Nothing was omitted from any retained span, so the package carries no omission record.  No retained line contains a citation, so the wrapper carries no bibliography and no bibliography entry.  No retained line contains a figure environment, an image-inclusion command or drawing code, so no image is delivered and the package declares no asset dependency.  Editorial formatting only, in addition: an AMS \texttt{amsart} preamble that restates the article's own declarations and macros used by the retained lines, namely the environments \texttt{theorem} on separate counters, together with the standard packages those lines need.  The retained environments print in this document as Theorem 1 in order of appearance, whereas the article prints the same items with the numbering of its own full document.  The complete native source of the article as distributed in the arXiv e-print for this exact version is bundled unchanged as \texttt{sources/196128/source0.tex}.
\begin{theorem}
Let $p$ be a prime, and let $K,L \subseteq \{0,1,\ldots,p-1\}$ be disjoint with $|K|=r$ and
$|L|=s$. Assume that $k>s-r\geq 0$ for every $k\in K$. If $F$ is a family of subsets of $[n]$ such that
\begin{enumerate}
\item[(i)] $|F| \in K+p\mathbb Z$ for every $F\in F$;
\item[(ii)] $|E\cap F| \in L+p\mathbb Z$ for all distinct $E,F\in F$,
\end{enumerate}
then
$$
|F|+\sum_{j=s-r+1}^{s}|N_j(F)| \le N(n,s,r).
$$
Equivalently,
$$
|F| \le \sum_{j=s-r+1}^{s} |\partial_j F|.
$$
\end{theorem}

\begin{proof}
Write $F=\{A_1,\ldots,A_m\}$ and order the sets so that
$$
|A_1|\le |A_2|\le \cdots \le |A_m|.
$$
For each $i$, define
$$
f_i(x):=\prod_{\ell\in L:\,\ell<|A_i|}(v_{A_i}\cdot x-\ell),
$$
where
$$
v_{A_i}\cdot x=\sum_{j=1}^n (v_{A_i})_j x_j.
$$

We claim that $f_i(A_i)\neq 0$ and $f_i(A_j)=0$ for all $j<i$. Indeed, if $\ell\in L$ and
$\ell<|A_i|$, then $|A_i|-\ell\neq 0$ in $\mathbb F_p$ because $|A_i| \pmod p \in K$ and
$K\cap L=\varnothing$. Hence $f_i(A_i)\neq 0$.

Now let $j<i$. Then $|A_j|\le |A_i|$ and $A_i\neq A_j$, so
$$
|A_i\cap A_j|<|A_i|.
$$
By (ii), there exists $\ell^*\in L$ such that
$$
|A_i\cap A_j|\equiv \ell^* \pmod p.
$$
Writing
$$
|A_i\cap A_j|=qp+\ell^*
$$
with $0\le \ell^*\le p-1$, we have
$$
\ell^*\le |A_i\cap A_j|<|A_i|.
$$
Therefore the factor indexed by $\ell^*$ appears in the product defining $f_i$, and it vanishes at
$A_j$. Thus $f_i(A_j)=0$ for every $j<i$. It follows that $f_1,\ldots,f_m$ are linearly
independent.

Define
$$
g(x):=\prod_{k\in K}\left(\sum_{j=1}^n x_j-k\right).
$$
If $A\in F$, then $|A|\pmod p\in K$, so $g(A)=0$. If $I\subseteq [n]$ satisfies $|I|\le s-r$, then
$|I|<k$ for every $k\in K$, hence $|I|-k\neq 0$ in $\mathbb F_p$ for every $k\in K$, and so
$g(I)\neq 0$. Lemma 3 therefore implies that
$$
\{x_Ig:\ |I|\le s-r\}
$$
is linearly independent over $\mathbb F_p$.

Let
$$
\mathcal J:=\bigcup_{j=s-r+1}^{s} N_j(F),
$$
and set
$$
B:=\{f_i:\ 1\le i\le m\}\cup \{x_Ig:\ |I|\le s-r\}\cup \{x_T:\ T\in\mathcal J\}.
$$
We claim that $B$ is linearly independent. Suppose that
$$
\sum_{i=1}^m \lambda_i f_i+\sum_{|I|\le s-r}\mu_Ix_Ig+\sum_{T\in\mathcal J}\nu_Tx_T=0.
$$
We first show that all $\lambda_i$ vanish. If some $\lambda_i\neq 0$, let $i_0$ be the smallest such
index. Evaluating at $A_{i_0}$, the middle sum vanishes because $g(A_{i_0})=0$, and the last sum
vanishes because $x_T(A_{i_0})=0$ for every $T\in\mathcal J$, since $T$ is contained in no member
of $F$. Also, $f_i(A_{i_0})=0$ for all $i>i_0$ by the triangular property proved above. Hence
$$
\lambda_{i_0}f_{i_0}(A_{i_0})=0,
$$
contradicting $f_{i_0}(A_{i_0})\neq 0$. Thus $\lambda_i=0$ for all $i$.

The remaining relation
$$
\sum_{|I|\le s-r}\mu_Ix_Ig+\sum_{T\in\mathcal J}\nu_Tx_T=0
$$
together with the facts that $|T|>s-r$ for every $T\in\mathcal J$ and that $g(I)\neq 0$ for all
$|I|\le s-r$ allows us to apply Lemma 4 with $t=s-r$ and $f=g$. Therefore all $\mu_I$ and
$\nu_T$ vanish, and $B$ is linearly independent.

Every element of $B$ is a multilinear polynomial of degree at most $s$: each $f_i$ has degree at
most $s$, each $x_Ig$ has degree $|I|+r\le s$, and each $x_T$ has degree $|T|\le s$. Hence
$$
|B|\le \dim\operatorname{span}\{x_I:\ |I|\le s\}
 =\sum_{i=0}^{s}\binom{n}{i}.
$$
On the other hand,
$$
|B|=|F|+\sum_{i=0}^{s-r}\binom{n}{i}+\sum_{j=s-r+1}^{s}|N_j(F)|.
$$
Subtracting $\sum_{i=0}^{s-r}\binom{n}{i}$ from both sides and using
$$
N(n,s,r)=\sum_{i=s-r+1}^{s}\binom{n}{i}
$$
gives
$$
|F|+\sum_{j=s-r+1}^{s}|N_j(F)|\le N(n,s,r).
$$
The equivalent form follows from
$$
|N_j(F)|=\binom{n}{j}-|\partial_jF|.
$$
\end{proof}

\end{document}
